\documentclass[12pt,a4paper]{article}

\usepackage[utf8]{inputenc}
\usepackage[T1]{fontenc}
\usepackage[english]{babel}
\usepackage{geometry}
\usepackage{graphicx}
\usepackage{hyperref}
\usepackage{listings}
\usepackage{xcolor}
\usepackage{booktabs}
\usepackage{parskip}

\geometry{margin=2.5cm}

\hypersetup{
    colorlinks=true,
    linkcolor=blue,
    urlcolor=blue
}

\lstset{
    basicstyle=\ttfamily\small,
    backgroundcolor=\color{gray!10},
    frame=single,
    breaklines=true,
    keywordstyle=\color{blue},
    commentstyle=\color{green!60!black},
    stringstyle=\color{orange},
    language=Python
}

% ---------------------------------------------------------------
\title{\textbf{Secure Messaging}\\
       \large RSA + AES Hybrid Encryption over HTTP}
\author{Denis O.}
\date{\today}

\begin{document}

\maketitle


% ---------------------------------------------------------------
\section{Objective}

Design and implement a secure messaging system between two clients in which messages are exchanged encrypted with an AES session key, itself distributed securely via RSA by a trusted key server.

% ---------------------------------------------------------------
\section{Architecture}

\subsection{Components}

The system comprises three processes:

\begin{itemize}
  \item \textbf{server.py} — key distribution server, listening on port~5000.
  \item \textbf{client.py (Alice)} — chat client, listening on port~5001.
  \item \textbf{client.py (Bob)} — chat client, listening on port~5002.
\end{itemize}

\subsection{High-level diagram}

\begin{center}
\begin{tabular}{ccc}
\toprule
\textbf{Alice (5001)} & \textbf{Server (5000)} & \textbf{Bob (5002)} \\
\midrule
\multicolumn{3}{c}{\texttt{POST /register} (public key) \(\rightarrow\) server} \\
\multicolumn{3}{c}{\texttt{POST /register} (public key) \(\rightarrow\) server} \\
\multicolumn{3}{c}{Alice: \texttt{POST /session(from=Alice, to=Bob)}} \\
\multicolumn{3}{c}{Server generates AES key; encrypts for Alice \& Bob} \\
\multicolumn{3}{c}{Alice decrypts her copy; forwards Bob's copy via \texttt{/receive\_key}} \\
\multicolumn{3}{c}{Alice \(\leftrightarrow\) Bob: AES-encrypted messages (\texttt{/receive\_message})} \\
\bottomrule
\end{tabular}
\end{center}

% ---------------------------------------------------------------
\section{Cryptographic Design}

\subsection{Key generation}

Each client generates a 2048-bit RSA key pair at startup:

\begin{lstlisting}
my_key     = RSA.generate(2048)
my_pub_key = my_key.publickey().export_key().decode('utf-8')
rsa_cipher = PKCS1_OAEP.new(my_key)
\end{lstlisting}

\subsection{Session key distribution}

The server generates a fresh 16-byte AES key per session and encrypts it separately for each party under their RSA public key (OAEP padding):

\begin{lstlisting}
session_key   = get_random_bytes(16)
enc_key_init  = PKCS1_OAEP.new(init_pub_key).encrypt(session_key)
enc_key_target = PKCS1_OAEP.new(target_pub_key).encrypt(session_key)
\end{lstlisting}

\subsection{Message encryption (AES-CBC)}

Messages are encrypted with AES-128 in CBC mode. The IV is prepended to the ciphertext and the whole blob is base64-encoded for JSON transport:

\begin{lstlisting}
def encrypt_aes(msg_text, key):
    cipher   = AES.new(key, AES.MODE_CBC)
    ct_bytes = cipher.encrypt(pad(msg_text.encode(), AES.block_size))
    return base64.b64encode(cipher.iv + ct_bytes).decode()

def decrypt_aes(b64_ct, key):
    raw    = base64.b64decode(b64_ct)
    iv, ct = raw[:16], raw[16:]
    cipher = AES.new(key, AES.MODE_CBC, iv)
    return unpad(cipher.decrypt(ct), AES.block_size).decode()
\end{lstlisting}

% ---------------------------------------------------------------
\section{API Reference}

\subsection{Server endpoints}

\begin{tabular}{lll}
\toprule
Method & Route & Description \\
\midrule
POST & \texttt{/register} & Store client username + RSA public key \\
POST & \texttt{/session}  & Generate \& distribute encrypted AES session key \\
\bottomrule
\end{tabular}

\subsection{Client endpoints}

\begin{tabular}{lll}
\toprule
Method & Route & Description \\
\midrule
GET  & \texttt{/}                 & Serve the browser chat UI \\
POST & \texttt{/send}             & Encrypt and send a message (triggers session setup if needed) \\
POST & \texttt{/receive\_key}     & Receive and decrypt the RSA-wrapped AES key \\
POST & \texttt{/receive\_message} & Receive, decrypt, and store an AES message \\
GET  & \texttt{/messages}         & Return chat history as JSON \\
\bottomrule
\end{tabular}

% ---------------------------------------------------------------
\section{How to Run}

\begin{lstlisting}[language=bash]
# Terminal 1
python server.py

# Terminal 2
python client.py 5001 Alice

# Terminal 3
python client.py 5002 Bob
\end{lstlisting}

Open \url{http://127.0.0.1:5001} (Alice) and \url{http://127.0.0.1:5002} (Bob) in two browser tabs. The UI auto-configures the counterpart target on load.

% ---------------------------------------------------------------
\section{Design Notes \& Limitations}

\begin{itemize}
  \item \textbf{Trust model}: the key server sees the AES session key in plaintext at generation time. A zero-knowledge improvement would replace this step with a Diffie–Hellman exchange.
  \item \textbf{No message authentication}: messages are encrypted but not signed. Adding RSA signatures (or HMAC) would prevent a MITM from injecting ciphertext.
  \item \textbf{Single session per client pair}: \texttt{current\_session\_key} is a single global variable; supporting multiple concurrent sessions would require a keyed dictionary.
  \item \textbf{No TLS}: all HTTP traffic is in cleartext at the transport layer. In production, every endpoint should be served over HTTPS.
\end{itemize}

% ---------------------------------------------------------------
\section{Possible Extensions}

\begin{itemize}
  \item RSA message signatures for authentication.
  \item Multi-session support (multiple concurrent chat pairs).
  \item Message timestamps and session-start notifications.
  \item Persistent local chat history.
\end{itemize}

% ---------------------------------------------------------------
\section{References}

\begin{itemize}
  \item PKCS \#1 OAEP — \href{https://datatracker.ietf.org/doc/html/rfc8017}{RFC 8017}
  \item PyCryptodome — \url{https://pycryptodome.readthedocs.io}
  \item Flask — \url{https://flask.palletsprojects.com}
\end{itemize}

\end{document}
